%% Based on a TeXnicCenter-Template by Gyorgy SZEIDL.
%%%%%%%%%%%%%%%%%%%%%%%%%%%%%%%%%%%%%%%%%%%%%%%%%%%%%%%%%%%%%

%------------------------------------------------------------
%
\documentclass[a4paper]{article}%
\usepackage{amsmath}%
\usepackage{amsfonts}%
\usepackage{amssymb}%
\usepackage{graphicx}
\usepackage[spanish]{babel}
\usepackage[latin1]{inputenc}
\usepackage{makeidx}
%-------------------------------------------
\newtheorem{theorem}{Theorem}
\newtheorem{acknowledgement}[theorem]{Acknowledgement}
\newtheorem{algorithm}[theorem]{Algorithm}
\newtheorem{axiom}[theorem]{Axiom}
\newtheorem{case}[theorem]{Case}
\newtheorem{claim}[theorem]{Claim}
\newtheorem{conclusion}[theorem]{Conclusion}
\newtheorem{condition}[theorem]{Condition}
\newtheorem{conjecture}[theorem]{Conjecture}
\newtheorem{corollary}[theorem]{Corollary}
\newtheorem{criterion}[theorem]{Criterion}
\newtheorem{definition}[theorem]{Definition}
\newtheorem{example}[theorem]{Example}
\newtheorem{exercise}[theorem]{Exercise}
\newtheorem{lemma}[theorem]{Lemma}
\newtheorem{notation}[theorem]{Notation}
\newtheorem{problem}[theorem]{Problem}
\newtheorem{proposition}[theorem]{Proposition}
\newtheorem{remark}[theorem]{Remark}
\newtheorem{solution}[theorem]{Solution}
\newtheorem{summary}[theorem]{Summary}
\newenvironment{proof}[1][Proof]{\textbf{#1.} }{\ \rule{0.5em}{0.5em}}

\begin{document}
\thispagestyle{empty}

\begin {center}

\medskip
UNIVERSIDAD DE BUENOS AIRES

Facultad de Ciencias Exactas y Naturales

Departamento de Computaci�n
\vspace{3cm}


\textbf{\Large Nuevas fuentes de informaci�n para entrenamiento de etiquetadores gramaticales}

\vspace{2cm}



\vspace{2cm}

\noindent Tesista: Fernando Jorge Rodriguez\\
\noindent Director: Dr. Jos� Casta�o


\vspace{1cm}
Buenos Aires, Marzo de 2012.
\vspace{1cm}
%\begin{center}

\includegraphics[scale=.6]{img/dc.jpg}
\end {center}

\newpage 
\tableofcontents
\newpage 
\subsection{Motivaci�n}
El etiquetado o anotado gramatical, tambi�n conocido como Part-of-speech tagging, POS tagging o simplemente POST, es el proceso de asignar una etiqueta gramatical a cada una de las palabras de un texto seg�n su categor�a l�xica. Como se menciona m�s adelante, el etiquetado gramatical juega un papel importante en  �reas de la linguistica computacional como por ejemplo s�ntesis del habla, reconocimiento del habla y recuperaci�n de la informaci�n.

Este proceso es realizado manualmente por linguistas o autom�ticamente mediante programas conocidos como etiquetadores gramaticales. Estos programas son entrenados con el objetivo de que aprendan a etiquetar. Este entrenamiento se realiza utilizando un corpus anotado previamente denominado corpus de entrenamiento. De esta manera el etiquetador obtiene, procesa y retiene informaci�n sobre como est� anotado el corpus de entrenamiento que utiliza posteriormente durante el proceso de etiquetaci�n gramatical. 

Una vez finalizado el entrenamiento, se procede a etiquetar el texto requerido utilizando la informacion adquirida.

Basicamente un corpus etiquetado o anotado es una lista de palabras con su correspondiente etiqueta gramatical. Por ejemplo:
\\ \\
A			DT\\
form			NN\\
of			IN\\
asbestos		NN\\
once			RB\\
used			VBN\\
to			TO\\
make			VB\\
Kent			NNP\\
cigarette		NN\\
filters			NNS\\
has			VBZ\\
caused			VBN\\
a			DT\\
high			JJ\\
percentage		NN\\
of			IN\\
cancer			NN\\
deaths			NNS\\
among			IN\\
a			DT\\
group			NN\\
of			IN\\
workers			NNS\\
exposed			VBN\\
to			TO\\
it			PRP\\
more			RBR\\
than			IN\\
30			CD\\
years			NNS\\
ago			RB\\
,			,\\
researchers		NNS\\
reported		VBD\\
.			.\\ \\

La complejidad del etiquetado gramatical reside en que la etiqueta no depende solo de la palabra o tipo de palabra que se est� etiquetando, por el contrario, la ubicaci�n y el contexto de la palabra en donde �sta aparece es un factor determinante al asignar la etiqueta correspondiente.

Como fu� mencionado anteriormente, la etapa de entrenamiento de un etiquetador gramatical obtiene, procesa y retiene informaci�n sobre el contexto y la ubicaci�n de una palabra y la etiqueta asignada. Luego se utiliza esta informaci�n en el proceso de etiquetacion, donde el etiquetador analiza para cada palabra su ubicaci�n y contexto, y en base a ello y al conocimiento adquirido previamente con el corpus de entrenamiento determina una etiqueta gramatical.

Uno de los grandes problemas del etiquetado gramatical reside en la falta de corpus anotados para utilizar durante el entrenamiento.

Los corpus etiquetados (tambi�n llamados corpus de entrenamiento) son etiquetados manualmente por linguistas especializados. Es un trabajo profundamente meticuloso y tedioso ya que el linguista debe dar una etiqueta gramatical palabra por palabra, en corpus del orden del MILLON de palabras. Ademas de lo tedioso y complejo del trabajo, el tiempo empleado para etiquetar un corpus es sumamente extenso y como consecuencia el valor economico es realmente alto, ya que intervienen grupos de trabajo altamente capacitados durante un tiempo prolongado.

El resultado de este complejo proceso artesanal es una tabla de palabras con su correspondiente etiqueta gramatical como se mostr� anteriormente. Ante la importancia que adquieren los corpus etiquetados es inevitable pensar en alg�n otro tipo de texto que posea informaci�n de etiquetas. 

Por ejemplo un diccionario contiene una palabra, su definici�n y algunos ejemplos en donde �sta aparece con cada uno de sus sentidos. Es decir que de alguna manera un diccionario contiene por cada palabra uno o m�s contextos en donde �sta aparece etiquetada. Entonces si tomamos todos los ejemplos de cada palabra de un diccionario y su etiqueta podemos construir un corpus parcialmente anotado. 

\subsection{Trabajo realizado}
La idea de este trabajo es suplir la falta de corpus de enternamiento utilizando la informaci�n de etiquetado que posee un diccionario como una nueva fuente de informaci�n para entrenar etiquetadores autom�ticos. Utillizar esta informaci�n existente para etiquetar, medir y comparar resultados en el rendimiento conseguido. Combinar e integrar esta informaci�n obtenida con corpus de entrenamiento cl�sicos e inspeccionar resultados en el etiquetado.
Este trabajo menciona detalladamente la forma de extraer la informaci�n relevante de etiquetas gramaticales a partir de un diccionario y las decisiones que fueron aplicadas.
Por �ltimo se realizan mediciones entre esta nueva fuente de informaci�n y los corpus cl�sicos de entrenamiento y se presentan las conclusiones.





















\section{The Most Important Features}

\noindent Sectioning commands. The first one is the\\
\hspace*{\fill} \verb"\section{The Most Important Features}" \hspace*{\fill}\\
command. Below you shall find examples for further sectioning commands:

\subsection{Subsection}
Subsection text.

\subsubsection{Subsubsection}
Subsubsection text.

\paragraph{Paragraph}
Paragraph text.

\subparagraph{Subparagraph}Subparagraph text.\vspace{2mm}

Select a part of the text then click on the button Emphasize (H!), or Bold (Fs), or
Italic (Kt), or Slanted (Kt) to typeset \emph{Emphasize}, \textbf{Bold},
\textit{Italics}, \textsl{Slanted} texts.

You can also typeset \textrm{Roman}, \textsf{Sans Serif}, \textsc{Small Caps}, and
\texttt{Typewriter} texts.

You can also apply the special, mathematics only commands $\mathbb{BLACKBOARD}$
$\mathbb{BOLD}$, $\mathcal{CALLIGRAPHIC}$, and $\mathfrak{fraktur}$. Note that
blackboard bold and calligraphic are correct only when applied to uppercase letters A
through Z.

You can apply the size tags -- Format menu, Font size submenu -- {\tiny tiny},
{\scriptsize scriptsize}, {\footnotesize footnotesize}, {\small small}, {\normalsize
normalsize}, {\large large}, {\Large Large}, {\LARGE LARGE}, {\huge huge} and {\Huge
Huge}.

You can use the \verb"\begin{quote} etc. \end{quote}" environment for typesetting
short quotations. Select the text then click on Insert, Quotations, Short Quotations:

\begin{quote}
The buck stops here. \emph{Harry Truman}

Ask not what your country can do for you; ask what you can do for your
country. \emph{John F Kennedy}

I am not a crook. \emph{Richard Nixon}

I did not have sexual relations with that woman, Miss Lewinsky. \emph{Bill Clinton}
\end{quote}

The Quotation environment is used for quotations of more than one paragraph. Following
is the beginning of \emph{The Jungle Books} by Rudyard Kipling. (You should select
the text first then click on Insert, Quotations, Quotation):

\begin{quotation}
It was seven o'clock of a very warm evening in the Seeonee Hills when Father Wolf woke
up from his day's rest, scratched himself, yawned  and spread out his paws one after
the other to get rid of sleepy feeling in their tips. Mother Wolf lay with her big gray
nose dropped across her four tumbling, squealing cubs, and the moon shone into the
mouth of the cave where they all lived. ``\emph{Augrh}'' said Father Wolf, ``it is time
to hunt again.'' And he was going to spring down hill when a little shadow with a bushy
tail crossed the threshold and whined: ``Good luck go with you, O Chief of the Wolves;
and good luck and strong white teeth go with the noble children, that they may never
forget the hungry in this world.''

\ jackal---Tabaqui the Dish-licker---and the wolves of India despise Tabaqui
because he runs about making mischief, and telling tales, and eating rags and pieces of
leather from the village rubbish-heaps. But they are afraid of him too, because
Tabaqui, more than any one else in the jungle, is apt to go mad, and then he forgets
that he was afraid of anyone, and runs through the forest biting everything in his way.
\end{quotation}

Use the Verbatim environment if you want \LaTeX\ to preserve spacing, perhaps when
including a fragment from a program such as:
\begin{verbatim}
#include <iostream>         // < > is used for standard libraries.
void main(void)             // ''main'' method always called first.
{
 cout << ''This is a message.'';
                            // Send to output stream.
}
\end{verbatim}
(After selecting the text click on Insert, Code Environments, Code.)


\subsection{Mathematics and Text}

It holds \cite{KarelRektorys} the following
\begin{theorem}
(The Currant minimax principle.) Let $T$ be completely continuous selfadjoint operator
in a Hilbert space $H$. Let $n$ be an arbitrary integer and let $u_1,\ldots,u_{n-1}$ be
an arbitrary system of $n-1$ linearly independent elements of $H$. Denote
\begin{equation}
\max_{\substack{v\in H, v\neq
0\\(v,u_1)=0,\ldots,(v,u_n)=0}}\frac{(Tv,v)}{(v,v)}=m(u_1,\ldots, u_{n-1})
\label{eqn10}
\end{equation}
Then the $n$-th eigenvalue of $T$ is equal to the minimum of these maxima, when
minimizing over all linearly independent systems $u_1,\ldots u_{n-1}$ in $H$,
\begin{equation}
\mu_n = \min_{\substack{u_1,\ldots, u_{n-1}\in H}} m(u_1,\ldots, u_{n-1}) \label{eqn20}
\end{equation}
\end{theorem}
The above equations are automatically numbered as equation (\ref{eqn10}) and
(\ref{eqn20}).

\subsection{List Environments}

You can create numbered, bulleted, and description lists using the tag popup
at the bottom left of the screen.

\begin{enumerate}
\item List item 1

\item List item 2

\begin{enumerate}
\item A list item under a list item.

The typeset style for this level is different than the screen style. \ The
screen shows a lower case alphabetic character followed by a period while the
typeset style uses a lower case alphabetic character surrounded by parentheses.

\item Just another list item under a list item.

\begin{enumerate}
\item Third level list item under a list item.

\begin{enumerate}
\item Fourth and final level of list items allowed.
\end{enumerate}
\end{enumerate}
\end{enumerate}
\end{enumerate}

\begin{itemize}
\item Bullet item 1

\item Bullet item 2

\begin{itemize}
\item Second level bullet item.

\begin{itemize}
\item Third level bullet item.

\begin{itemize}
\item Fourth (and final) level bullet item.
\end{itemize}
\end{itemize}
\end{itemize}
\end{itemize}

\begin{description}
\item[Description List] Each description list item has a term followed by the
description of that term. Double click the term box to enter the term, or to
change it.

\item[Bunyip] Mythical beast of Australian Aboriginal legends.
\end{description}

\subsection{Theorem-like Environments}

The following theorem-like environments (in alphabetical order) are available
in this style.

\begin{acknowledgement}
This is an acknowledgement
\end{acknowledgement}

\begin{algorithm}
This is an algorithm
\end{algorithm}

\begin{axiom}
This is an axiom
\end{axiom}

\begin{case}
This is a case
\end{case}

\begin{claim}
This is a claim
\end{claim}

\begin{conclusion}
This is a conclusion
\end{conclusion}

\begin{condition}
This is a condition
\end{condition}

\begin{conjecture}
This is a conjecture
\end{conjecture}

\begin{corollary}
This is a corollary
\end{corollary}

\begin{criterion}
This is a criterion
\end{criterion}

\begin{definition}
This is a definition
\end{definition}

\begin{example}
This is an example
\end{example}

\begin{exercise}
This is an exercise
\end{exercise}

\begin{lemma}
This is a lemma
\end{lemma}

\begin{proof}
This is the proof of the lemma.
\end{proof}

\begin{notation}
This is notation
\end{notation}

\begin{problem}
This is a problem
\end{problem}

\begin{proposition}
This is a proposition
\end{proposition}

\begin{remark}
This is a remark
\end{remark}

\begin{solution}
This is a solution
\end{solution}

\begin{summary}
This is a summary
\end{summary}

\begin{theorem}
This is a theorem
\end{theorem}

\begin{proof}
[Proof of the Main Theorem]This is the proof.
\end{proof}
\medskip

This text is a sample for a short bibliography. You can cite a book by making use of
the command \verb"\cite{KarelRektorys}": \cite{KarelRektorys}. Papers can be cited
similarly: \cite{Bertoti97}. If you want multiple citations to appear in a single set
of square brackets you must type all of the citation keys inside a single citation,
separating each with a comma. Here is an example: \cite{Bertoti97, Szeidl2001,
Carlson67}.

\begin{thebibliography}{9}                                                                                                %
\bibitem {KarelRektorys}Rektorys, K., \textit{Variational methods in Mathematics,
Science and Engineering}, D. Reidel Publishing Company,
Dordrecht-Hollanf/Boston-U.S.A., 2th edition, 1975

\bibitem {Bertoti97} \textsc{Bert\'{o}ti, E.}:\ \textit{On mixed variational formulation
of linear elasticity using nonsymmetric stresses and displacements}, International
Journal for Numerical Methods in Engineering., \textbf{42}, (1997), 561-578.

\bibitem {Szeidl2001} \textsc{Szeidl, G.}:\ \textit{Boundary integral equations for
plane problems in terms of stress functions of order one}, Journal of Computational and
Applied Mechanics, \textbf{2}(2), (2001), 237-261.

\bibitem {Carlson67}  \textsc{Carlson D. E.}:\ \textit{On G\"{u}nther's stress functions
for couple stresses}, Quart. Appl. Math., \textbf{25}, (1967), 139-146.
\end{thebibliography}


\appendix

\section{The First Appendix}

The appendix fragment is used only once. Subsequent appendices can be created
using the Section Section/Body Tag.
\end{document}
